\section{A planar proof for a fixed number of petals}\label{sec:plane}
Proposition~\ref{prop:comparison} asks for a placement with two properties.
Sharing a word among all \(r\) groups must cost more than sharing it
within the groups separately. At the same time, checking the agreement
rules must have bounded cost in each coordinate, however large its
alphabet. We will place two groups on line segments and the other
\(r-2\) groups at points between them. In this placement, comparing
two symbols already shared over suitable groups of sites costs a bounded
amount.

\subsection{Comparing and copying shared symbols}
Fix a spectral point \(\varphi\) on \(S\). Let \(P,T\subseteq S\) be
disjoint nonempty sets of sites and put \(K=P\cup T\). Each site of \(K\)
receives a \(b\)-bit control. An \emph{equality filter} consists of
fixed nonzero input and return tensors \(R,Q\) and controlled local maps.
When the sites of \(P\) all read \(x\) and those of \(T\) all read \(y\),
the maps must take \(R\) to \(\mathbf 1_{\{x=y\}}Q\). On any other
assignment of the controls the output is unrestricted. Denote by
\(c_b(P,T)\) the infimum of the net costs \(h(R)-h(Q)\).
Taking \(b=\lceil\log_2|\Sigma|\rceil\) accommodates any finite alphabet
\(\Sigma\). A one-symbol alphabet needs no test.

\par\Needspace{6\baselineskip}
\begin{theorem}[Equality testing]\label{thm:filter}
Put \(g=w(P)+w(T)-w(K)\), \(g_-=\max\{0,-g\}\), and
\(\ell=\log_2(2|K|)\). For every \(b\ge1\),
\[
        g_- b\ \le\ c_b(P,T)\
             \le\ g_- b+Cw(K)\sqrt{b\ell}.
\]
In particular the cost per bit tends to \(g_-\). When \(g>0\), there
is an equality filter of net cost at most \(Cw(K)^2\ell/g\) for every
\(b\), where \(C\) is an absolute constant.
\end{theorem}

For an implication \((U,V)\to W\), let \(P\) read sites in both \(U,V\)
and let \(T\) read sites in \(W\). When the implication fails, \(P\)
reads one symbol and \(T\) a different one, so the filter gives zero.
When all three symbols agree it returns \(Q\). Thus a positive \(g\)
gives exactly the alphabet-independent guarantee of
Proposition~\ref{prop:comparison}. Sections~\ref{sec:copy} and~\ref{sec:filter}
prove the equality estimate. We will also need a second way to enforce
an implication when the target lies between the two sources.

\begin{lemma}[Copying inside a convex hull]\label{lem:segment}
Let \(P,Z\) be disjoint nonempty sets of sites with \(w(P\cup Z)=w(P)\).
For every \(N\ge1\), a controlled restriction on \(P\) returns
\(Q\otimes E_Z(i)\) whenever all sites of \(P\) read \(i\in[N]\),
for a fixed nonzero tensor \(Q\). Its net cost is at most
\(w(P)\log_2(|Z|+2)\). For the Euclidean spectral point of
Theorem~\ref{thm:flat}, the width condition holds whenever
\(Z\subseteq\conv(P)\).
\end{lemma}
\begin{proof}
Write \(c=|Z|\). Choose \(L\) from \(c+1,c+2\) so the characteristic
does not divide \(L\), and let \(\zeta\) be a primitive \(L\)-th root
of unity. Supply \(G_{P\cup Z}(L^N)\), indexing the shared basis by
\(t=(t_1,\ldots,t_N)\in(\mathbb Z/L\mathbb Z)^N\). When a source
reads \(i\), it deletes \(t_i\) and keeps the other coordinates in
order. One source also multiplies by \(L^{-1}\zeta^{-ct_i}\).
Each target maps \(e_t\) to \(\sum_{j=1}^N\zeta^{t_j}e_j\).

Fix the retained source index and a tuple of target indices
\((j_z)_{z\in Z}\). If \(m\) of them equal \(i\), summation over \(t_i\)
contributes the factor \(L^{-1}\sum_{s=0}^{L-1}\zeta^{(m-c)s}\).
This is zero unless \(m=c\), since \(0\le m\le c<L\).
When \(m=c\) all phases cancel, leaving coefficient one. The output is
therefore \(G_P(L^{N-1})\otimes E_Z(i)\). With \(w(P\cup Z)=w(P)\),
Equation~\eqref{eq:price} gives net cost \(w(P)\log_2 L\).
\end{proof}

\subsection{The placement and its agreement rules}
Name the two anchor groups \(A,B\) and the remaining groups
\(Z_1,\ldots,Z_{r-2}\). In each coordinate use the rules
\begin{equation}\label{eq:rules}
 (A,B)\to Z_i,\qquad (A,Z_i)\to B,\qquad (B,Z_i)\to A,
 \qquad (Z_i,Z_j)\to A\quad(i<j).
\end{equation}
If the anchors agree, the first rules force agreement everywhere. An
agreeing anchor and satellite force the other anchor to agree. Two
agreeing satellites first force agreement with \(A\), which in turn
forces agreement with \(B\). Thus a column satisfying the rules is
either constant or pairwise distinct.

Set
\[
   t=\frac{1}{2(r-1)},\qquad \eta=\frac{t^2}{4},\qquad
   y_i=1-\frac{i}{r-1}\quad (1\le i\le r-2).
\]
The anchors occupy the ends of two vertical unit segments at horizontal
coordinates \(0\) and \(1+\eta\), together with the sites
\(a_i=(0,y_i)\) and \(b_i=(1+\eta,y_i)\) on those segments.
Group \(Z_i\) occupies the single site \((t,y_i)\). The whole placement
has semiperimeter \(D=2+\eta\). By Theorem~\ref{thm:flat}, a group
\(P\) has weight \(w(P)=\delta(P)/D\), with \(\delta(P)\) its
semiperimeter. Only the anchors contribute to the sum of group weights,
so the surplus for Proposition~\ref{prop:comparison} is
\[
              \rho=1-\frac{2}{2+\eta}=\frac{\eta}{2+\eta}>0 .
\]

\begin{figure}[htb]
\centering
\begin{tikzpicture}[x=1cm,y=1cm,font=\small]
  \begin{scope}
    \draw[gray!65] (0,0) rectangle (3.9,2.65);
    \fill[rust!13] (0,0)--(.65,1.8)--(0,2.65)--cycle;
    \draw[rust!75!black] (0,0)--(.65,1.8)--(0,2.65);
    \draw[very thick,ink] (0,0)--(0,2.65) (3.9,0)--(3.9,2.65);
    \draw[sea,thick,dashed] (0,1.8)--(3.9,1.8);
    \foreach \x in {0,3.9}{\foreach \y in {0,.65,1.8,2.65}
      {\fill[ink] (\x,\y) circle(1.6pt);}}
    \node[above] at (0,2.68) {\(A\)};
    \node[above] at (3.9,2.68) {\(B\)};
    \node[left] at (0,1.8) {\(a_i\)};
    \node[right] at (3.9,1.8) {\(b_i\)};
    \fill[rust] (.65,1.8) circle(2.2pt) node[above right,text=black] {\(Z_i\)};
    \fill[rust] (.65,.65) circle(2.2pt) node[right,text=black] {\(Z_j\)};
    \draw[<->,gray!80!black] (0,-.32)--node[below,text=black]{\(1+\eta\)}(3.9,-.32);
  \end{scope}
  \begin{scope}[xshift=6.3cm]
    \fill[rust!10] (0,.2) rectangle (1.4,2.35);
    \draw[gray!75] (0,.2) rectangle (1.4,2.35);
    \draw[very thick,ink] (0,.2)--(0,2.35);
    \draw[very thick,rust] (1.4,.2)--(1.4,2.35);
    \foreach \y/\lab in {2.35/i,.2/j}{
      \fill[ink] (0,\y) circle(1.8pt) node[left,text=black]{\(a_{\lab}\)};
      \fill[rust] (1.4,\y) circle(2.2pt) node[right,text=black]{\(Z_{\lab}\)};}
    \draw[<->,gray!80!black] (-.72,.2)--node[left,text=black]{\(d\)}(-.72,2.35);
    \draw[<->,gray!80!black] (0,-.32)--node[below,text=black]{\(t\)}(1.4,-.32);
  \end{scope}
\end{tikzpicture}
\caption{At left, \(a_i\) and \(b_i\) can copy a common symbol to
\(Z_i\) on their joining segment. The shaded triangle is the convex hull of \(A\cup\{Z_i\}\),
whose semiperimeter exceeds that of \(A\).
At right, the four sites used to check \((Z_i,Z_j)\to A\)
form a rectangle with \(d=|y_i-y_j|>t\).}
\label{fig:placement}
\end{figure}

For \((A,B)\to Z_i\), read \(a_i\) and \(b_i\). Their segment contains
\(Z_i\), as in Figure~\ref{fig:placement}. Lemma~\ref{lem:segment}
copies their common label to \(Z_i\) at bounded cost. The local map
there accepts the copied label precisely when it equals the satellite's
own control.

For \((Z_i,Z_j)\to A\), use \(P=\{Z_i,Z_j\}\) and \(T=\{a_i,a_j\}\).
Both segments have length \(d=|y_i-y_j|\ge2t\), and their union spans
the rectangle in Figure~\ref{fig:placement}. The corresponding
semiperimeter gap is \(2d-(d+t)=d-t\ge t\). Theorem~\ref{thm:filter}
applies with \(g=(d-t)/D\). For \(r=3\) there is only one satellite
and no rule of this type.

For \((A,Z_i)\to B\), let \(P\) contain \(Z_i\) and the endpoints of
\(A\), and let \(T\) consist of the endpoints of \(B\). Adding a point
at height \(y\in[0,1]\) and horizontal distance \(s\) from a unit
segment increases its semiperimeter by
\[
  \frac{\sqrt{y^2+s^2}+\sqrt{(1-y)^2+s^2}-1}{2}
     \ \ge\ \frac{\sqrt{1+4s^2}-1}{2} .
\]
The inequality follows from the triangle inequality applied to
\((y,s)\) and \((1-y,s)\). Here \(s=t\). For \((B,Z_i)\to A\) the
same calculation uses \(s=1+\eta-t\ge t\). In either case the
increase is at least
\(2t^2/(\sqrt{1+4t^2}+1)>t^2/2\), since \(t\le1/4\).
The two anchors together span the whole rectangle. Subtracting its
semiperimeter \(2+\eta\) leaves a gap greater than \(t^2/2-\eta=\eta\).
Thus both filters have \(g>\eta/D\).

Each filter just used reads at most five sites. The \(O(r)\) rules
with a satellite and an anchor as antecedents therefore cost \(O(r^2)\)
each, and the \(O(r^2)\) satellite-pair rules cost \(O(r)\) each.
Copying from the two anchors costs \(O(1)\) per satellite.
Since \(\rho\) is of order \(r^{-2}\), Proposition~\ref{prop:comparison}
gives \(\log_2|\W|=O(r^5 k)\). The transfer in Section~\ref{sec:words}
adds only \(O(k)\) to this logarithm for set families.
This already gives an exponential bound in \(k\) for every fixed \(r\).
Part~\ref{part:quant} obtains the stronger dependence on \(r\) by
testing many implications together.
